\documentclass[11pt]{article}
\usepackage[a4paper,margin=25mm]{geometry}
\usepackage{amsmath,amssymb,mathtools}
\usepackage{fontspec,unicode-math}
\setmainfont{Latin Modern Roman}
\setsansfont{Latin Modern Sans}
\setmathfont{Latin Modern Math}
\usepackage{xurl,hyperref}
\hypersetup{colorlinks=true,linkcolor=blue,urlcolor=blue,pdftitle={Problem statement}}
\newcommand{\sref}[2]{\hyperlink{source:#1:#2}{[S#2]}}
\newcommand{\cataloguescope}[1]{\paragraph{Mathematical question.}#1}
\begin{document}
\section{Lebesgue's universal covering problem on the minimum-area convex shape in the plane that can cover any shape of diameter one}\label{statement}

\subsection{Definitions and mathematical statement}

\paragraph{Shared notation.}$\mathbb N=\{1,2,\ldots\}$, and $\mathbb Z,\mathbb Q,\mathbb R,\mathbb C$ denote the integers, rationals, reals and complex numbers. Any other conventions are specified locally.


For $A\subset\mathbb R^2$, its diameter is $\sup\{|x-y|:x,y\in A\}$. A convex universal cover is a convex planar set $K$ containing an isometric copy of every set $A$ of diameter one. An isometry may translate, rotate or reflect $A$. Area means two-dimensional Lebesgue measure.
\paragraph{Statement recorded in the supplied source.}
Find a convex universal cover of minimum area and determine that minimum. \sref{statement}{2}
\subsection{Short English statement}
What is the smallest area of a convex planar set that can contain, after a rigid motion, every planar set of diameter one?
\subsection{Sources}
\begin{description}
\item[\hypertarget{source:statement:1}{S1}] ULAM.AI and the UnsolvedMath Contributors, UnsolvedMath: A Curated Collection of Open Mathematics Problems, catalogue entry. CC BY 4.0. \url{https://huggingface.co/datasets/ulamai/UnsolvedMath}.
\item[\hypertarget{source:statement:2}{S2}] Philip Gibbs, A New Slant on Lebesgue’s Universal Covering Problem (2014); arXiv:1401.8217. \url{https://arxiv.org/pdf/1401.8217}.
\end{description}
\label{end:statement}
\end{document}
